\chapter{11-9-26}
\section{La relación de divisibilidad}

\begin{definicion}
Sean $a,b \in \mathbb{Z}$. Decimos que $a$ \textbf{divide} a $b$, y
escribimos $a \mid b$, si existe un entero $c$ tal que $b=ac$.
\end{definicion}

\begin{definicion}
Fijado un entero $m$, decimos que $a$ es \textbf{congruente} con $b$
módulo $m$, y escribimos $a \equiv b \pmod m$, si $m \mid a-b$.
\end{definicion}

\begin{nota}
La relación $a \mid b$ es reflexiva y transitiva, pero no es
antisimétrica en $\mathbb{Z}$: si $a\mid b$ y $b \mid a$, entonces
$a=b$ o $a=-b$ (no necesariamente $a=b$).

En cambio, la relación de congruencia módulo $m$ \textbf{sí} es una
relación de equivalencia en $\mathbb{Z}$ (es reflexiva, simétrica y
transitiva).
\end{nota}

\begin{teoria}
Al fijar el módulo $m$, la relación de congruencia parte a
$\mathbb{Z}$ en \textbf{clases de equivalencia}. El conjunto de todas
estas clases se anota
\[
\mathbb{Z}_m = \{ [a] : a \in \mathbb{Z} \},
\]
donde $[a] = \{ b \in \mathbb{Z} : b \equiv a \pmod m \}$ es la clase
de equivalencia de $a$.
\end{teoria}

\begin{ejemplo}
\textbf{Caso $m=1$.} Como $1$ divide a cualquier entero, $a \equiv b
\pmod 1$ vale para todo $a,b \in \mathbb{Z}$. Por lo tanto hay una
única clase:
\[
\mathbb{Z}_1 = \{ [0] \}.
\]

\textbf{Caso $m=0$.} Se tiene $a\equiv b \pmod 0$ si y sólo si
$0 \mid a-b$, es decir, existe $c\in\mathbb{Z}$ tal que $a-b = 0\cdot c
= 0$, o sea $a=b$. Así que cada entero forma su propia clase.

\textbf{Caso $m=2$.} $a \equiv b \pmod 2$ si y sólo si $2 \mid a-b$,
es decir $a-b$ es par. Hay exactamente dos clases:
\[
[0] = \{\text{números pares}\}, \qquad [1] = \{\text{números impares}\},
\]
\[
\mathbb{Z}_2 = \{[0],[1]\}.
\]

\textbf{Caso $m=3$.} $a \equiv b \pmod 3$ si y sólo si $3 \mid a-b$,
es decir $a-b$ es múltiplo de $3$, o equivalentemente $a=b+3c$ para
algún $c\in\mathbb{Z}$. Entonces
\[
[a] = \{ a+3c : c \in \mathbb{Z} \}.
\]
En particular,
\[
[0] = \{3c : c\in\mathbb{Z}\}, \qquad [1] = \{1+3c : c\in\mathbb{Z}\}.
\]
Estas dos clases son disjuntas, pero no cubren todo $\mathbb{Z}$: por
ejemplo el $2$ no está en ninguna de las dos (si $2=3c$ entonces
$c=2/3 \notin \mathbb{Z}$; si $2=1+3c$ entonces $c=1/3\notin
\mathbb{Z}$). Agregamos entonces
\[
[2] = \{2+3c : c\in\mathbb{Z}\}.
\]
Notar que $[3]=[0]$, $[4]=[1]$, $[5]=[2]$, y en general las clases se
repiten cada $3$ enteros. Así,
\[
\mathbb{Z}_3 = \{[0],[1],[2]\}.
\]
\end{ejemplo}

\begin{proposicion}
Para todo $n \in \mathbb{N}_0$, $[n] \in \{[0],[1],[2]\}$ (clases
módulo $3$).
\end{proposicion}

\begin{demostracion}
Por inducción en $n$.

\textbf{Caso base.} Para $n=0,1,2$ la afirmación es inmediata, pues
$[0],[1],[2]$ son justamente los elementos del conjunto
$\{[0],[1],[2]\}$.

\textbf{Paso inductivo.} Supongamos que $[n]\in\{[0],[1],[2]\}$ para
cierto $n\in\mathbb{N}_0$ (hipótesis inductiva). Entonces existe un
entero $c$ tal que $n=3c$, $n=3c+1$ o $n=3c+2$.

\begin{enumerate}
    \item Si $n=3c$, entonces $n+1=3c+1$, así que $[n+1]=[1]$.
    \item Si $n=3c+1$, entonces $n+1=3c+2$, así que $[n+1]=[2]$.
    \item Si $n=3c+2$, entonces $n+1=3c+3=3(c+1)$, así que
    $[n+1]=[0]$.
\end{enumerate}

En cualquiera de los tres casos, $[n+1]\in\{[0],[1],[2]\}$. Por
inducción, la afirmación vale para todo $n\in\mathbb{N}_0$.
\end{demostracion}

\begin{ejemplo}
Mostrar que para todo $n \in \mathbb{N}_0$ vale que $[-n] \in
\{[0],[1],[2]\}$ (clases módulo $3$).
\end{ejemplo}

\begin{demostracion}
Por inducción en $n$.

\textbf{Caso base.} Para $n=0$, $[-0]=[0]\in\{[0],[1],[2]\}$.

\textbf{Paso inductivo.} Supongamos que $[-n]\in\{[0],[1],[2]\}$.
Veamos que $[-(n+1)]\in\{[0],[1],[2]\}$, usando que
$-(n+1)=-n-1$.

\begin{enumerate}
    \item Si $[-n]=[0]$, entonces $-n=3c$ para algún $c$, y
    $-(n+1) = 3c-1 = 3(c-1)+2$, así que $[-(n+1)]=[2]$.
    \item Si $[-n]=[1]$, entonces $-n=3c+1$, y
    $-(n+1) = 3c$, así que $[-(n+1)]=[0]$.
    \item Si $[-n]=[2]$, entonces $-n=3c+2$, y
    $-(n+1) = 3c+1$, así que $[-(n+1)]=[1]$.
\end{enumerate}

En cualquiera de los tres casos, $[-(n+1)]\in\{[0],[1],[2]\}$. Por
inducción, la afirmación vale para todo $n\in\mathbb{N}_0$.
\end{demostracion}

\begin{nota}
Combinando las dos proposiciones anteriores, todo entero $n\in
\mathbb{Z}$ (sea positivo, negativo o cero) satisface $[n]\in
\{[0],[1],[2]\}$, así que efectivamente $\mathbb{Z}_3=\{[0],[1],[2]\}$.
De la misma manera, $\mathbb{Z}_2 = \{[0],[1]\}$.
\end{nota}

\begin{proposicion}
Si $m$ es un entero positivo, entonces
\[
\mathbb{Z}_m = \{[0],[1],\dots,[m-1]\}
\]
y tiene exactamente $m$ elementos (todos distintos entre sí).
\end{proposicion}

Para demostrar esta proposición en general vamos a necesitar el
\textbf{algoritmo de la división}.

\begin{proposicion}[{title={Proposición: algoritmo de la división}}]
Sea $m$ un entero positivo. Si $a$ es un entero, entonces existen
enteros $q$ y $r$ tales que
\[
a = qm+r, \qquad 0 \le r < m,
\]
y $q,r$ están unívocamente determinados por $a$ y $m$.
\end{proposicion}

\begin{demostracion}
\textbf{Existencia, caso $a \ge 0$.}
Consideremos el conjunto
\[
Q = \{ b \in \mathbb{N}_0 : a < bm \}.
\]
Como $m\ge 1$, se tiene $a < (a+1)m$, así que $a+1 \in Q$ y $Q \neq
\varnothing$. Por el principio de buena ordenación, $Q$ tiene un
mínimo $b_0$.

Como $b_0$ es el mínimo de $Q$, se cumple $a<b_0 m$ pero $b_0-1
\notin Q$ (notar $b_0 \ge 1$, ya que si $b_0=0$ sería $a<0$,
imposible), es decir
\[
a \ge (b_0-1)m.
\]
Juntando ambas desigualdades:
\[
(b_0-1)m \le a < b_0 m.
\]
Definimos $q=b_0-1$ y $r=a-qm$. Entonces $a=qm+r$ y, de la
desigualdad anterior,
\[
0 \le r < m.
\]

\textbf{Existencia, caso $a<0$.}
Sea $a'=-a>0$. Por el caso anterior existen $q_1,r_1\in\mathbb{Z}$
tales que
\[
a' = q_1 m + r_1, \qquad 0\le r_1<m.
\]
Entonces $a=-a' = -q_1 m - r_1$.

\begin{itemize}
    \item Si $r_1=0$: tomamos $q=-q_1$, $r=0$, y $a=qm+r$ con
    $0\le r<m$.
    \item Si $r_1>0$: escribimos
    \[
    a = -q_1 m - r_1 = (-q_1-1)m + (m-r_1).
    \]
    Como $0<r_1<m$, se tiene $0<m-r_1<m$. Tomamos $q=-q_1-1$,
    $r=m-r_1$, y de nuevo $a=qm+r$ con $0\le r<m$.
\end{itemize}

\textbf{Unicidad.} Supongamos que $q,r$ y $q',r'$ son enteros tales
que
\[
a=qm+r, \quad 0\le r<m, \qquad a=q'm+r', \quad 0\le r'<m.
\]
Restando ambas igualdades:
\[
r-r' = (q'-q)m.
\]
Como $0\le r<m$ y $0\le r'<m$, se tiene $-m<r-r'<m$, así que
\[
-m < (q'-q)m < m.
\]
Dividiendo por $m>0$: $-1<q'-q<1$, y como $q'-q$ es entero, debe ser
$q'-q=0$, es decir $q=q'$. Luego, de $r-r'=(q'-q)m=0$ se obtiene
$r=r'$.
\end{demostracion}

\begin{demostracion}[{title={Demostración de la proposición $\mathbb{Z}_m=\{[0],\dots,[m-1]\}$}}]
Sea $x \in \mathbb{Z}_m$. Entonces $x=[a]$ para algún $a\in\mathbb{Z}$.
Por el algoritmo de la división, existen $q,r\in\mathbb{Z}$ con
$a=qm+r$ y $0\le r<m$. Como $a-r=qm$, se tiene $a\equiv r \pmod m$,
así que
\[
x = [a] = [r] \in \{[0],[1],\dots,[m-1]\}.
\]
Esto prueba que $\mathbb{Z}_m \subseteq \{[0],\dots,[m-1]\}$, y la
inclusión recíproca es trivial, así que
$\mathbb{Z}_m=\{[0],\dots,[m-1]\}$.

Falta ver que estas $m$ clases son todas distintas. Sean
$r,r'\in\{0,1,\dots,m-1\}$ con $[r]=[r']$. Entonces $r\equiv r'
\pmod m$, es decir $m \mid r-r'$. Como $0\le r,r'<m$, se tiene
$|r-r'|<m$, y el único múltiplo de $m$ en ese rango es el $0$; por lo
tanto $r-r'=0$, o sea $r=r'$.

En conclusión, las clases $[0],[1],\dots,[m-1]$ son distintas dos a
dos, así que $\mathbb{Z}_m$ tiene exactamente $m$ elementos.
\end{demostracion}

\section{Compatibilidad de la congruencia con las operaciones}

\begin{proposicion}
Sean $a,a',b,b' \in \mathbb{Z}$ tales que $a \equiv a' \pmod m$ y
$b \equiv b' \pmod m$. Entonces
\[
a+b \equiv a'+b' \pmod m \qquad \text{y} \qquad ab \equiv a'b' \pmod m.
\]
\end{proposicion}

\begin{demostracion}
Por hipótesis existen enteros $c,d$ tales que
\[
a-a' = cm, \qquad b-b' = dm.
\]

\textbf{Suma.}
\[
(a+b)-(a'+b') = (a-a')+(b-b') = cm+dm = (c+d)m,
\]
así que $m \mid (a+b)-(a'+b')$, es decir $a+b\equiv a'+b' \pmod m$.

\textbf{Producto.}
\[
ab-a'b' = ab-a'b+a'b-a'b' = (a-a')b + a'(b-b') = cmb + a'dm =
m(cb+a'd),
\]
así que $m \mid ab-a'b'$, es decir $ab\equiv a'b' \pmod m$.
\end{demostracion}

\begin{ejemplo}
Calcular el resto de dividir $3^{10}$ por $10$.

Como $3^2=9\equiv -1 \pmod{10}$, usando la compatibilidad del
producto:
\[
3^4 = (3^2)^2 \equiv (-1)^2 = 1 \pmod{10},
\]
\[
3^8 = (3^4)^2 \equiv 1^2 = 1 \pmod{10}.
\]
Entonces
\[
3^{10} = 3^8 \cdot 3^2 \equiv 1\cdot(-1) = -1 \equiv 9 \pmod{10}.
\]
Por lo tanto, el resto de dividir $3^{10}$ por $10$ es $9$.
\end{ejemplo}